\section{An absolute bound for the local orbital operator}
\label{sec:absolute}
\hypertarget{source-absolute-amplitude}{}

The squared norm controls the cell. A nonzero preceding minor permits an additional pointwise estimate.

Let \(\Delta_r\) denote the leading minor of order \(r\) of the matrix of pairings between the two frames. A Gram label with nonzero Gram invariant and \(\Delta_r\ne 0\) gives a normal form for the orbit of two flags. We use its leading minors to control the
coordinates of that form and estimate the orbital operator
with the normalization from Section~\ref{sec:local-amplitudes}.
The bounds in terms of the Gram determinant and the nonzero
minor sizes are used in the global row sum of
Section~\ref{sec:row-sums}, which uses an \(L^2\) bound
when \(\Delta_r=0\). For the two frames of Definition~\ref{hyp:paired-frame-moment-range}, the shell bound is \eqref{eq:side-moment}; the volume on each shell is measured in the free-root quotient of \eqref{eq:fibre-volume}.

Fix a real place and two words. Work on a cell with nonzero Gram
invariant and use
\[
 D=|\Delta_r|>0,\quad m=2r+2,\quad n=2r+k+1,\quad h=k-r-2,
 \qquad r\ge1,\quad k\ge2r+6.
\]
Let \(D_e\) be the maximum of one and the reciprocals of the
absolute values of the nonzero leading minors. In a block antidiagonal representative the
nonzero leading minors of positive order occur at block ends.
Taking a proper initial segment of the next block gives a
singular submatrix. Thus \(D_e\) is the same maximum that enters
the height estimates in Section~\ref{sec:row-sums}.

Use the Witt space \(X\oplus W\oplus X^*\) with
\(\dim X=r\), \(\dim W=k+1\) and \(W\) of signature \((k,1)\).
Take as frame an ordered basis of \(X\) followed by the vector
\(z\) on the negative line of \(W\). Write \(f_0,\ldots,f_r\) for this frame and \(F_{a-1}\) for the
span of the columns before \(a\). Define
\[
 l_a(x)=\operatorname{dist}(x^{-1}f_a,x^{-1}F_{a-1})
 \quad(0\le a\le r),\qquad S(x)=\max(1,\max_a l_a(x)).
\]
The corresponding primed functions use the primed frame ending in \(z'\).

Let \(\Phi(x)=B(\sigma(x)v_0)\) and \(\Phi'(x)=B'(\sigma(x)v_0)\),
where \(v_0\) belongs to the lowest compact type. Both Bessel
maps use the same discrete series. The quotient measures for the
two Bessel subgroups are denoted by \(d\nu_S,d\nu_{S'}\).
If the definite representations have dimensions \(d_W,d_{W'}\),
the normalizing scalar \(c_N>0\) satisfies
\[
 \int\|\Phi\|^2d\nu_S=c_Nd_W,\qquad
 \int\|\Phi'\|^2d\nu_{S'}=c_Nd_{W'}.
\]
For the right-side stabilizer
\(R=N'\cap\gamma^{-1}N\gamma\) let \(d\nu_R\) be the quotient
of \(dg\) by its specified Haar measure. The operator is
\[
 J=c_N^{-1}\int_{R\backslash G}
       \Phi(\gamma x)\otimes\Phi'(x)^*\,d\nu_R(x).
\]

\subsection{Length moments for both frames}

Two vectors occur in the integrand, in different frames. Moment bounds for their lengths restrict the region that matters before its volume is estimated.

At the standard character parameter write \(B_0\) for the
Pieri weight bound and \(B_v=b_1\) for the smallest absolute value in the negative tail at this real place.
For a unit highest-weight vector in the lowest compact type let
\(dP\) be the squared normalized Bessel coefficient times the
quotient measure. It has mass one in a nonzero model. Let
\(T_z\) be the positive norm parameter of the negative line.
The two
estimates needed from the coefficient recursion are
\[
 \int l_0^{2j}dP\le A^j,\qquad
 \int(l_{a+1}/l_a)^{2j}dP\le B_0^{4j},
 \qquad A=n^2 T_z B_0^4(6B_v)^{2m}.
\]

\begin{paramdefinition}[Uniform range for the frame moments]\label{hyp:paired-frame-moment-range}
At a marked real place, use the ordered row and the interlacing
definite column of Definition~\ref{hyp:critical-level-weights}.
Write \(\lambda\) for the row, \(u\) for the column, and put
\[
  B_v=b_1,\quad B_0=b^{\max(2,C_u)},\quad
  \gamma=nm^{10},\quad q_0=km^5 .
\]
The ranges for the row, column and intermediate integral shifts are
\begin{align}
 0\le\lambda_a\le B_v/8\ (a<k),\qquad
 &-2B_v\le\lambda_a\le-B_v\ (a\ge k),\qquad
 \|\lambda\|_\infty,\|u\|_\infty\le B_0,\label{eq:marked-pieri-windows}\\
 \lambda_a-\lambda_{a+1}>\gamma,\qquad
 &q_0\le\gamma/8,\qquad 100q_0,\ 10^6n^7\le B_v,\qquad
 \|\nu\|_1\le q_0 .\nonumber
\end{align}
The initial parameter has \(\|\nu\|_1\le C_nkm^3\).
Use \(m\ge |C_n|+25\), \(k\ge m^4\), \(2v\le b\),
\(1\le g\le b^{2C_g}\), and \(T_z\le b^{C_Tm}\) in both orientations.
\end{paramdefinition}

\begin{auxlemma}[The marked pair is in the common moment range]\label{lem:marked-paired-moment-range}
For the marked pair of words at one real place, let \(j\) be the minor order, \(q\) the chosen moment order and \(T_z\) the norm of the distinguished negative line in either orientation. Uniformly on the selected family,
\[
    \log T_z\ll m\log b,\qquad n^2,\,6B_v\le B_0^2,
    \qquad \|\nu\|_1\le qj\ll km^3 .
\]
In particular Definition~\ref{hyp:paired-frame-moment-range} applies at both words.
\end{auxlemma}
\begin{proof}
For the selected pair, the marked negative norm is \(T_z=T_u\) or
\(T_{u'}\) in the two orientations, so Proposition~\ref{prop:assigned-prime} gives
\[
    0<\log T_z\ll m\log b .
\]
Both type primes are at most \(b^{C_g}\), and
\(2v=\operatorname{rd}(E)\le b\). To check \eqref{eq:marked-pieri-windows}, denote the central spacing
of Lemma~\ref{hyp:height-section-weights} by \(h=\eta_{\rm size}/m\).
Its proof gives, with indices now starting at one,
\[
 b_m-b_1\le (m-1)(h+1)\le 2\eta_{\rm size}\le b_1 ,
 \qquad b_1\ge(t-1)\eta_{\rm size}.
\]

The same window selection in Paragraph~\ref{par:window-successive-bounds}
allows the following polynomial lower thresholds:
\[
 \min(4\Delta,h)> C_\ell k m^4+\gamma,\qquad
 b_1\ge\max(32k\Delta,k^{40}) .
\]
Indeed the first fixes lower bounds for \(\Delta\) and \(h\);
increasing \(\eta_{\rm size}\) by a factor polynomial in \(x\)
after choosing \(\Delta\) gives the second. The logarithmic
upper threshold for the width is twice that for \(\Delta\), as in
the proof of Lemma~\ref{hyp:height-section-weights}.
Definition~\ref{hyp:critical-level-weights} gives the head coordinates
in \([0,4k\Delta]\) with consecutive spacing \(4\Delta\).
The gap at the junction is at least \(b_1\). Hence the unshifted
bounds hold for both words, with the same row. For the definite column the upper endpoint is \(u_1\le b^{C_u}\).
Interlacing and the bound on the row give
\[
  -b^2\le\lambda_k<u_k\le u_i\le u_1\le b^{C_u},\qquad
  \|u\|_\infty\le B_0 .
\]

Since \(n=k+m-1\le k^2\) and \(k\ge1000\) at the common threshold,
\[
 10^6 n^7\le k^{16}\le b_1,\qquad
 100q_0=100km^5\le k^{8}\le b_1,\qquad
 8q_0\le nm^{10}.
\]
In particular unit shifts of total size at most \(q_0\) preserve
the ordering and leave gaps at least \(\gamma-q_0\ge2\).
The half-integral displacement between the definite and ambient
lists and their interlacing follow from the binary choices in
Definition~\ref{hyp:critical-level-weights}.
Also \(B_v\le B_0\), so \(6B_v\le6B_0\le B_0^2\).
The moment choice in
Lemma~\ref{lem:row-moment-budget} permits at most \(qj\) unit shifts
for a minor of order \(j\). The explicit comparison in the row sum is
\[
   \|\nu\|_1\le qj\le C_{\rm pr}km^3,\qquad
   k\ge m^5(\log b)x^2.
\]
It gives the length and rank conditions simultaneously for both words.
\end{proof}

The tensor with the dual standard representation carries an invariant
Hermitian form of mixed signature. To estimate the norms of its
projections we first separate the diagonal and off-diagonal parts of
the split Casimir. Throughout this calculation put
\[
  \mu=\lambda+\nu,\quad
  I=\{1,\ldots,k-1\},\quad I'=\{k,\ldots,n\},\quad |I'|=m .
\]

Write \(E_{ab}\) for the complex matrix units acting on the
discrete series, and \(\varepsilon_a\) for their coordinate signs.
The compact roots join coordinates of the same sign. Let
\(\rho_c,\rho_n\) be the half sums of the positive compact and
noncompact roots, respectively, in the order defined by \(\mu\).
Thus the lowest compact type has highest weight \(L\), and its
compact Casimir eigenvalue \(K\) and the full eigenvalue \(c\) are
\begin{equation}\label{eq:paired-lowest-casimirs}
 \rho=\rho_c+\rho_n,\quad L=\mu+\rho_n-\rho_c,\quad
 K=\|\mu+\rho_n\|_2^2-\|\rho_c\|_2^2,\quad
 c=\|\mu\|_2^2-\|\rho\|_2^2 .
\end{equation}
Say that a vector has level at most \(j\) if its compact types
have compact Casimir at most \((\sqrt K+4nj)^2\), and every
diagonal weight \(\omega\) in that vector satisfies
\(\|\omega-L\|_1\le j\). The highest vector has level zero.

\begin{auxlemma}[Matrix units on the permitted compact types]\label{lem:paired-lowest-casimir}
For the weights of Lemma~\ref{lem:marked-paired-moment-range}, at
any level \(j\le4n^2+2\),
\[
  n\sum_{a\ne b}\|E_{ab}w\|^2\le (B_v/16)^2\|w\|^2,
  \qquad \sum_{a,b}\|E_{ab}w\|^2\le B_0^4\|w\|^2 .
\]
\end{auxlemma}
\begin{proof}
The adjoint relation \(E_{ab}^*=\varepsilon_a\varepsilon_b E_{ba}\)
expresses the sum of the squares as twice the compact Casimir minus
the full Casimir. For each occurring diagonal weight one has
\[
 \sum_a\omega_a^2\ge \sum_a L_a^2-2\|L\|_\infty j
 \ge\sum_a L_a^2-8B_vj .
\]
Indeed \(|\mu_a|\le 3B_v\) and
\(|L_a-\mu_a|\le n\) by the root counts.

Orthogonality of diagonal weight spaces now bounds the off-diagonal
square sum by \(D_j\|w\|^2\), where
\[
 D_j=2(\sqrt K+4nj)^2-c-\sum_a L_a^2+8B_vj .
\]

The terms of order \(B_v^2\) cancel. Coordinatewise the unshifted
summand of \(2K-c-\sum L_a^2\) is, with
\(U=(\rho_n)_a,\ T=(\rho_c)_a\),
\[
       2\mu_a(U+T)+(U+T)^2-4T^2+\rho_a^2 .
\]
The root counts give \(|U|,|T|,|U+T|,|\rho_a|\le n\).
By \eqref{eq:paired-lowest-casimirs}, \(\sqrt K\le4nB_v\).
Using \(j\le5n^2\), we obtain
\begin{align*}
 D_j &\le n(12nB_v+4n^2)+64n^2j B_v+32n^2j^2+8B_vj\\
     &\le 500n^4B_v \le \frac{B_v^2}{256n}.
\end{align*}

The last two inequalities use \(10^6 n^7\le B_v\).
The unshifted coordinates have magnitude at most \(2B_v\);
the initial shift costs at most \(B_v/100\), and the lowest-type
shift and the level cost at most \(n+j\). Hence
\[
 \max_a |\omega_a|\le 2B_v+B_v/100+n+j\le3B_v,\qquad
 9nB_v^2+B_v^2/(256n)\le B_0^4
\]
since \(B_0\ge B_v\ge10^6n^7\). Adding the diagonal squares proves
the second inequality.
\end{proof}

\begin{auxlemma}[Sum of the squared dual Pieri projections]\label{lem:paired-spectral-projections}
In the dual standard tensor, write \(R_s\) for the projection onto
the unitary summand with parameter \(\mu-e_s\), endowed with the norm
induced up to sign by the invariant Hermitian form. For any coordinate
vector \(e_c\) and any \(w\) of level \(j\le4n^2\),
\[
        \sum_s\|R_s(e_c\otimes w)\|^2
                       \le (6B_v)^{2m}\|w\|^2 .
\]
\end{auxlemma}
\begin{proof}
Write \(x=e_c\otimes w\) and let \(p_s\) project back into the tensor
product. In the positive direct-sum norm, its invariant form and
split Casimir are
\[
 [y,y]=\sum_a\varepsilon_a\|y_a\|^2,\qquad
 (Sy)_a=-\sum_b E_{ab}y_b,\qquad
 Sp_s=s_s p_s,\quad s_s=-\mu_s+(1-n)/2 .
\]
These eigenspaces and their signed isometries are the dual Pieri
decomposition used in Lemma~\ref{lem:weighted-fitted-translate}.
The distinct eigenvalues differ by integers; the eigenspaces are
orthogonal for \([\,,\,]\) since \(S\) is self-adjoint for this form.

Work on the finite-dimensional span \(V_x=\sum_s\C p_s x\).
The coordinates of elements of \(V_x\) have level at most \(j+2\).
In fact the coordinate injection, equivariant projection and
coordinate extraction have combined weight change \(e_a-e_c\).

For the compact type, the coordinate maps factor through tensoring
by the standard module and its dual. In one such tensor step the
compact infinitesimal action is \(X\otimes 1+1\otimes X\).
Using a compact orthonormal basis \((X_i)\) gives
\[
 \left(\sum_i\|X_i y\|^2\right)^{1/2}
 \le (C+\sqrt n)\|y\|
\]
on a tensor of a type of Casimir at most \(C^2\) with a standard
or dual standard module. Indeed each input Casimir is bounded there
by \(C^2\) and \(n\), and the triangle inequality applies in the
direct sum over \(i\). Compact equivariance transports this bound
to the image compact types. 
In the abstract shifted summand, a weight \(\omega-e_c\) is
measured from \(L-e_s\), so
\[
 \|(\omega-e_c)-(L-e_s)\|_1\le j+2 .
\]

The corresponding \(\sqrt K\) changes by at most \(\sqrt n\). To see the latter estimate, the two compact
highest weights differ by a standard weight; each type occurs
in the tensor of the other with the standard or its dual by
the compact Pieri rule. Thus the bound \(4n(j+2)\) in the definition
of the new level allows both tensor operations. The same argument
applies to one matrix unit, since their compact module is a summand
of the standard tensor its dual.

Write \(S=D+O\) with \((Dy)_a=-E_{aa}y_a\). The two intervals
containing the spectrum of the individual diagonal entries on these
weight spaces, and also the numbers \(s_a\), are
\[
  [-5B_v/32,B_v/32]\quad(a\in I),\qquad
  [31B_v/32,65B_v/32]\quad(a\in I').
\]
Indeed the change from \(-\lambda_a\) to a diagonal eigenvalue is
at most \(B_v/100+n+j+2\le B_v/32\).
By Cauchy--Schwarz in each row and Lemma~\ref{lem:paired-lowest-casimir},
\[
 \|Oy\|^2\le n\sum_b\sum_{a\ne b}\|E_{ab}y_b\|^2
            \le (B_v/16)^2\|y\|^2,\qquad \|Sy\|\le3B_v\|y\|
 \quad(y\in V_x).
\]

On the circle \(|\zeta+B_v/16|=9B_v/16\), both intervals have
distance at least \(15B_v/32\) from \(\zeta\). Subtracting the bound
for \(O\) and inverting on \(V_x\) gives
\[
 \|(\zeta-S)^{-1}\|\le \frac{32}{13B_v},\qquad
 P=\sum_{s\in I}p_s=\frac{1}{2\pi i}\oint(\zeta-S)^{-1}\,d\zeta,
 \qquad \|P\|\le18/13 .
\]
For \(s\in I'\) interpolate only on the remaining \(m\) eigenvalues.
Their separation is at least one, and \(\|1-P\|<3\), so
\[
 p_s=\prod_{t\in I',\ t\ne s}\frac{S-s_t}{s_s-s_t}(1-P),
 \qquad \|p_s\|\le 3(6B_v)^{m-1}.
\]

For \(s\in I\) and \(y=p_s x\), the equation
\((D-s_s)y=-Oy\) on the coordinates \(I'\) has diagonal gap
at least \(30B_v/32\). Thus
\[
  \sum_{a\in I'}\|y_a\|^2\le \|y\|^2/225,\qquad
  [y,y]\ge (1-2/225)\|y\|^2 ,
\]
because every coordinate in \(I\) has positive sign.
The positive summands therefore add without cancellation in
\([Px,Px]\), whereas each remaining signed norm is bounded by its
positive direct-sum norm. It follows that
\[
 \begin{aligned}
 \sum_s\|R_sx\|^2
 &= [Px,Px]+\sum_{s\in I'}|[p_s x,p_s x]|\\
 &\le\big((18/13)^2+9m(6B_v)^{2(m-1)}\big)\|x\|^2
 \le(6B_v)^{2m}\|w\|^2 .
 \end{aligned}
\]
Here the final comparison uses \(m\le n\le B_v\).
\end{proof}

\begin{auxlemma}[A single dual-standard translation of a Bessel period]
\label{lem:paired-lowering-factor}
Let \(\lambda_1>\cdots>\lambda_n\) be a marked ambient
Harish--Chandra parameter, \(u_1>\cdots>u_k\) the compact
parameter, and \(\mu=\lambda+\nu\) with
\(\|\nu\|_1+1\le q_0\). Take a nonzero Bessel model at \(\mu\)
and its normalized period maps as in
Proposition~\ref{prop:real-period-norm}. Let \(\kappa_s\) be the scalar
obtained by inserting the distinguished covector in the dual standard
tensor, projecting to the discrete series of parameter \(\mu-e_s\)
and applying its period map. In the same Haar measure at both parameters,
\[
 |\kappa_s|^2\frac{d_\mu}{d_{\mu-e_s}}
   =T_z R_s\le T_z B_0^4,\qquad
 R_s=\frac{\prod_{i=1}^k|\mu_s-\tfrac12-u_i|}
              {\prod_{t\ne s}|\mu_s-1-\mu_t|}.
\]
Here \(d_\mu\) is the formal degree and \(-T_z\) the squared
Hermitian norm of the distinguished negative vector.
\end{auxlemma}

\begin{proof}
\proofstep{Compute the projected period.}
The shifted summands in the dual Pieri rule have multiplicity one.
On the \(s\)-summand the split Casimir has eigenvalue
\(-\mu_s+(1-n)/2\), by subtracting the two factor Casimirs from
the tensor Casimir. Let \(P_s\) denote its spectral projection. In the
Capelli column determinant use the central polynomial \(Q(Y)\)
and its signed comatrix \(C(Y)\), with argument \(Y+(n-1)/2\).
The quantum comatrix identity gives
\[
 Q(Y)=\prod_t(Y+\mu_t),\qquad
 C(Y)P_t=\prod_{j\ne t}(Y+\mu_j)P_t .
\]
Indeed multiplication by the shifted split Casimir in the comatrix
identity multiplies the \(t\)-summand by \(Y+\mu_t\); polynomial
division gives the second equality, including at the roots.
In particular \(C(-\mu_s)=\prod_{t\ne s}(\mu_t-\mu_s)P_s\).
The determinant and minor conventions are those of
\cite[Proposition~2.7(4) and Remarks~2.8(1), 2.11]{mno-yangians}.

Use the ordered Witt basis consisting of the two null flags and
the middle block. In the Bessel functional the comatrix entry
between the outer frame vectors reduces to the compact \(k\)-row
minor. To see the factors, number the basis positions from
\(0\) to \(2r+k\). For \(r>0\) the cofactor deletes the last row
and the first column. Expand first in the first remaining row,
then in the last remaining column, successively down the two flags.
The nonzero root character entries outside the middle are
\[
 (a,a+1)\quad(a+1<r\text{ or }a>r+k),\qquad
 (r-1,r+k),\quad(r+k,r+k+1).
\]

The ordinary flag entries have modulus one and the two joining
entries each have modulus \(\sqrt{T_z}\). At each outer expansion
the displayed entry is the only one on which the functional survives.
In a column expansion the argument drops by one, from commuting
the last matrix unit to the left using
\([E_{ij},E_{ab}]=\delta_{ja}E_{ib}-\delta_{ib}E_{aj}\).
After the \(r\) column removals the compact argument at the root is
\[
 -\mu_s+\tfrac{n-1}{2}-r=-\mu_s+\tfrac k2,\qquad
 Q_{\rm compact}(-\mu_s+\tfrac k2)
       =\prod_{i=1}^k(-\mu_s+\tfrac12+u_i).
\]
The latter equality is the compact Capelli central character. For
\(r=0\) the cofactor is already the compact block, and insertion
and contraction of the negative vector give \(T_z\) directly.

The tensor Hermitian form restricts on the summand to its unitary
form times a sign \(\epsilon_s\). Thus the period of the reverse
contraction is \(\epsilon_s\overline{\kappa_s}\) times the shifted
period. Indeed pairing the two contractions on smooth vectors in
the period integral moves the adjoint across the Hermitian form,
and the inserted covector is fixed by the period subgroup.
On a vector with nonzero original period, the composite therefore
has modulus \(|\kappa_s|^2\). The cofactor evaluation gives
\[
 |\kappa_s|^2
 =T_z\frac{\prod_i|\mu_s-\tfrac12-u_i|}
              {\prod_{t\ne s}|\mu_s-\mu_t|},\qquad
 \frac{d_\mu}{d_{\mu-e_s}}
 =\prod_{t\ne s}\frac{|\mu_s-\mu_t|}{|\mu_s-1-\mu_t|}.
\]
The second identity follows by cancelling unchanged factors in
Harish--Chandra's Vandermonde formula at fixed Haar measure.
The admissible shifts remain regular, so all the denominators
are nonzero. Multiplication gives the asserted equality.

\proofstep{Pair the factors before estimating the quotient.}
In Definition~\ref{hyp:critical-level-weights} the upper entries
are \(v_1,\ldots,v_{k-1}\) and the remaining entries are
\(-b_1,\ldots,-b_m\). Use \(B_v=b_1\), \(q_0=km^5\) and the
spacing bound \(\gamma=nm^{10}\) from
Lemma~\ref{lem:marked-paired-moment-range}. Its selected row bounds give
\[
 \begin{gathered}
  0\le\lambda_i\le B_v/8\ (i<k),\qquad
  -2B_v\le\lambda_i\le-B_v\ (i\ge k),\\
  |u_i|,|\lambda_t|\le B_0,\qquad
  q_0\le\min(\gamma/8,B_v/4,B_0).
 \end{gathered}
\]
Strict interlacing and opposite parity at the compact and ambient
parameters give \(\lambda_i+\tfrac12\le u_i\) for \(i\le k\)
and \(u_i+\tfrac12\le\lambda_{i-1}\) for \(2\le i\le k\).
Put \(N_i=|\mu_s-\tfrac12-u_i|\),
\(D_t=|\mu_s-1-\mu_t|\) and \(\delta_t=|\nu_t|\).
Each \(N_i\le4B_0\), while \(D_t\) for \(t\ne s\) is a positive
integer and hence at least one.

If \(s<k\), assign a distinct denominator to each indicated numerator:
\[
 \begin{array}{c|c|c}
   i&t&A\rule{0pt}{2.5ex}\\ \hline
   2\le i<s&i-1&\lambda_t+1-\mu_s\\
   s+2\le i\le k&i&\mu_s-1-\lambda_t
 \end{array}
\]
In the first line \(u_i\le\lambda_{i-1}-1/2\) and
\(u_i\ge\lambda_i+1/2>\mu_s-1/2\).
In the second, \(u_i\ge\lambda_i+1/2\) and
\(u_i\le\lambda_{i-1}-1/2\le\mu_s-1/2\).
The separation of \(\lambda_s\) from \(\lambda_i\) in the first
line and from \(\lambda_{i-1}\) in the second is at least
\(\gamma\), so in both lines
\[
 N_i\le A,\qquad D_t\ge A-\delta_t,\qquad
 A\ge\gamma/2,\qquad 2\delta_t\le A .
\]

If \(s\ge k\), pair every \(2\le i<k\) with \(t=i-1\), and put
\(A=\lambda_t+1-\mu_s\). Here \(\lambda_t,\lambda_i\ge0\) and
\(\mu_s\le-B_v+q_0\), so \(N_i\le A\), \(D_t\ge A-\delta_t\),
\(A\ge B_v/2\) and \(2\delta_t\le A\).
For each paired factor, taking \(c=\gamma/2\) or \(B_v/2\)
according to the case gives
\[
 \frac{N_i}{D_t}\le\frac{A}{A-\delta_t}
      \le 1+\frac{2\delta_t}{A}
      \le e^{2\delta_t/c},\qquad
 \prod_{i\,\mathrm{paired}}\frac{N_i}{D_{t(i)}}\le e^{2q_0/c}\le e .
\]
The upper case omits at most \(1,s,s+1\), the lower case at most
\(1,k\). Unused denominators are at least one. Consequently
\[
 R_s\le
 \begin{cases}64e B_0^3&s<k,\\16e B_0^2&s\ge k.\end{cases}
 \qquad R_s\le B_0^4\quad(B_0\ge64e).
\]
The last threshold holds uniformly at the selected family scale.
\end{proof}


\begin{proposition}[Length moments uniform in the two words]\label{prop:paired-frame-moments}
For either marked word of Lemma~\ref{lem:marked-paired-moment-range}
and an initial shift as in Definition~\ref{hyp:paired-frame-moment-range}, let \(l\) be any orthogonalized frame length and \(F\) its Bessel coefficient. In the measure of that moment, 
\[
 \int l^{2q}|F|^2\le H^{2q}\!\int|F|^2,\qquad q=24r^3+1.
\]
 The same \(H\) works for the primed and unprimed
frames.
\end{proposition}

\begin{proof}
\proofstep{The first length and successive ratios.}

For the first bound, tensor the coefficient with the dual standard
module and sum over an orthonormal basis of that module. The sum
of squares of the finite coefficients is \(l_0^2\). There are at
most \(n\) shifted parameters and \(n\) basis vectors. The
squared lowering coefficient, multiplied by the ratio of the two
formal degrees, is at most \(T_z B_0^4\) by
Lemma~\ref{lem:paired-lowering-factor}. For a given basis vector
Lemma~\ref{lem:paired-spectral-projections} bounds the sum of
squared projection norms by \((6B_v)^{2m}\) times the original squared norm. Thus the
single-step multiplier is at most
\[
           n^2(T_z B_0^4)(6B_v)^{2m}=A .
\]
An intermediate parameter may leave the range with a nonzero
Bessel model for this definite column. At that parameter the
period map and its contribution to the sum vanish. Starting from
each remaining nonzero coefficient, the same bound applies at
the next step. Iteration gives \(A^j\).

For the ratio, the simple root carrying the
orthogonalized \(a\)-th frame vector to the next has a Lie
algebra representative of Hilbert--Schmidt norm at most \(2l_a\).
Its character derivative has modulus \(2l_{a+1}\).
Differentiate the Bessel relation along this root. After dividing
its Cayley-basis coefficients by \(2l_a\), their squared moduli sum
to at most one. If \(F_v\) denotes the coefficient of the vector
\(v\), Cauchy--Schwarz gives
\[
   |F_v|^2(l_{a+1}/l_a)^2\le\sum_{i,j}|F_{E_{ij}v}|^2,\qquad
   \sum_{i,j}\|E_{ij}v\|^2\le B_0^4\|v\|^2 .
\]
Lemma~\ref{lem:paired-lowest-casimir} supplies the second
inequality. Each application
raises the level by two. These recursions use
\(\|\nu\|_1+j\le q_0\) and \(2j\le4n^2\) at the lowest type,
where \(\nu\) is the integral shift and \(q_0=km^5\).
Thus each intermediate vector lies in the range of the two lemmas.
\label{cor:absolute-amplitude-1}

\proofstep{An arbitrary frame length and lowest-type vector.}

Write \(l_a=l_0\prod_{i<a}(l_{i+1}/l_i)\). H\"older with
\(a+1\) equal exponents bounds its \(2q\)-moment by
\((A B_0^{4r})^q\), using the preceding estimates up to
\(j=(r+1)q\). The numerical hypotheses
\(n^2\le B_0^2\), \(6 B_v\le B_0^2\) give
\[
 A B_0^{4r}\le \max(1,T_z)B_0^{20(r+1)}.
\]
The weighted squared integral is a compact-invariant quadratic
form on the lowest type. Indeed right multiplication by the
compact group preserves each length and the integral on the
quotient. Finiteness on a highest vector implies finiteness on
its finite linear span of translates.

On an irreducible compact type, Schur's lemma makes the
form a scalar. More generally, on an isotypic component
\(V_\lambda\otimes M_\lambda\), its representing operator has the form
\[
     I_{V_\lambda}\otimes A_\lambda .
\]
Testing the highest-weight line tensored with the multiplicity
vectors bounds \(A_\lambda\) and hence the entire form. For the chosen compact type the
Bessel line is one dimensional. Replacing its distinguished map
by the actual map multiplies both integrals by the same factor. If \(F'=cF\), then
\[
  \frac{\int l_a^{2q}|F'|^2\,d\nu_S}{\int|F'|^2\,d\nu_S}
       =\frac{|c|^2\int l_a^{2q}|F|^2\,d\nu_S}{|c|^2\int|F|^2\,d\nu_S}
       \qquad(c\ne0).
\]

\proofstep{Transport and the common range.}

The moment estimate concerns one frame. The product integral requires the same moment order for the two transported distributions, even when the words differ.


For the actual character, conjugate by the diagonal torus
element taking \(2\pi/(vg)\) to the standard parameter, where
\(g=g_0g_1\). This changes \(l_a\) by a factor of modulus at
most \(\max(\pi/(vg),vg/\pi)^r\). It transports the quotient
measure in both integrals. With \(2v\le b\), \(1\le g\le
b^{2C_g}\), \(T_z\le b^{C_T m}\), and \(B_0=b^{\max(2,C_u)}\)
we obtain
\begin{equation}\label{eq:side-moment}
 \int l_a^{2q}\|\Phi\|^2d\nu_S
 \le H^{2q}\int\|\Phi\|^2d\nu_S.
\end{equation}
Here \(H=b^{m\max(C_M,0)}\) with
\(C_M=2C_g+2+|C_T|+10\max(2,C_u)\).
Interchanging the two words exchanges the negative lines,
definite representations, Bessel maps and quotient measures.

It preserves the common representation, the weight window, the
scalar \(c_N\), and the size hypotheses. Thus the same \(H\)
works for the primed moment. To check the uniform range, put
\(q=24r^3+1\le25r^3\). For
\(\|\nu\|_1\le C_nkm^3\), take \(m\ge |C_n|+25\)
and \(k\ge m^4\). The required interpolation and degree ranges follow from
\[
     (r+1)q\le mk,\qquad
     \|\nu\|_1+(r+1)q\le km^5,\qquad (r+1)q\le2n^2 .
\]
The spectral radius changes the threshold for \(m\).
The constant \(C_M\) is independent of that radius.

\end{proof}

\begin{construction}
\begin{tabular}{@{}ll@{}}\toprule
Factor of \(l_a=l_0\prod_{i<a}(l_{i+1}/l_i)\) & Bound for its \(2j\)-moment\\\midrule
Initial length \(l_0\) & \(A^j\)\\
Each of the \(a\) successive ratios & \(B_0^{4j}\)\\\midrule
Whole product, at exponent \(2q\) &
\(\displaystyle\int l_a^{2q}\,dP\le (A B_0^{4a})^q\)\\\bottomrule
\end{tabular}
\caption{Combining frame-length moments at the standard character. H\"older uses \(j=(a+1)q\) for the \(a+1\) factors, within the common range for the two words.}\label{fig:frame-length-moment-factors}
\end{construction}

\subsection{A volume estimate in the common stabilizer quotient}

Bounding the two lengths puts the pair of frames in a bounded region. The measure of that region must be taken modulo their common stabilizer, so counting all frame coordinates separately would give the wrong dimension.

Let \(E\) be the matrix of a basis of the isotropic flag and
\(Z\) its matrix of Hermitian pairings with the opposite
isotropic flag. Thus \(|\det Z|=D\). For the orientation
\((z,z',\gamma)\) let \(U_c\) be the Hermitian orthogonal
complement of the two full frames \((E,z')\) and
\(\gamma^{-1}(E,z)\). Project the columns of
\(x^{-1}\gamma^{-1}E\) Euclidean-orthogonally onto
\[
 P_x=\{p\mid H_{\rm Witt}p=p,\ xp\in U_c+X\}.
\]
Define \(Y\) to be their Euclidean wedge norm divided by
\(D\). Define \(Y'\) by reversing the frames and using
\(\gamma^{-1},\gamma x\). Below let \(S=S(\gamma x)\) and
\(S'=S'(x)\). \begin{lemma}[Gram and wedge bounds for two frames]\label{lem:paired-frame-wedges}
For \(D=|\det Z|>0\) and the wedge norms on \(P_x\) and its
frame-reversed counterpart, every \(x\) in the quotient chart satisfies
\[
 D\le S^r(S')^r,\qquad Y\le S^r/D,\qquad
 Y'\le(S')^r/D,\qquad YY'D\le2^r.
\]
If either projected frame has rank less than \(r\), its wedge vanishes.
\end{lemma}

\begin{proof}
\proofstep{Determinants and frame pairings.}

Apply successive column subtractions to each full flag frame,
with determinant one on its column indices. The Euclidean norm
of the wedge is the product of the orthogonalized lengths.
The two \(r\)-column wedges have norms at most \(S^r\) and
\((S')^r\). Pairing the wedges by the exterior power of the Euclidean
isometry \(H_{\rm Witt}\) and projecting by orthogonal contractions
gives
\[
     D\le S^r(S')^r,\qquad DY\le S^r,\qquad DY'\le(S')^r .
\]

For the last bound, let \(P,Q\) be the column maps of the two
projected frames. The operator obtained by contracting them
with \(Z^{-1}\) between the frame indices has norm at most
two. To see this, take vectors \(p,q\) in the two positive
projection spaces. In common coordinates their images under
\(\gamma x\) split as \(y_1+\xi_1,y_2+\xi_2\), with the
\(y_i\) perpendicular to both frames and \(\xi_1\in\gamma X\),
\(\xi_2\in X\). The mixed \(y,\xi\) terms vanish. Pairing
against the opposite flag recovers the coordinates of each
\(\xi_i\) by inverting \(Z\).

Write \([\ ,\ ]\) for the Witt form. Orthogonality gives the
contraction in absolute value as
\[
     |[\xi_1,\xi_2]|=|\langle q,p\rangle-[y_1,y_2]| ,
\]
with the conjugation convention immaterial to this equality of
absolute values. Its
restriction to the common complement is positive semidefinite.
Indeed orthogonality to \(X\) kills the dual isotropic coordinates
and orthogonality to the negative line leaves a positive
definite quotient by \(X\). The isotropic and mixed terms vanish on each diagonal,
so \([y_1,y_1]=\|p\|^2\), \([y_2,y_2]=\|q\|^2\).
The contraction has absolute value bounded by
\[
  |\langle q,p\rangle|+|[y_1,y_2]|\le 2\|p\|\|q\|.
\]
Its \(r\)-th exterior power therefore has norm at most
\(2^r\). This norm is \(\|\bigwedge^r P\|D^{-1}\|\bigwedge^r Q\|=YY'D\),
also when a column matrix is rank deficient.

\end{proof}

Write \(dN=\kappa_N\,dn\) in the unipotent matrix coordinates
and \(dR=s_R\,d\rho\) in the free roots of a normal form. The
following fibre bound explains the dependence on \(k\). At fixed
Iwasawa diagonal and middle Levi component, restrict to
\(S,S'\le R_0\), \(y<Y\le2y\), with \(R_0\ge1\), \(y>0\).
The Lebesgue volume in the unipotent coordinates modulo the
free roots is at most
\begin{equation}\label{eq:fibre-volume}
 (4r)^{11r^2}R_0^{20r^3}D_e^{12r^2}
       (\pi e/h)^{rh}(2y)^{2h}.
\end{equation}

The wedge bounds restrict the projected columns of the two frames. The remaining coordinates run over a fibre in the common stabilizer quotient and contribute its volume.

\begin{proposition}[Fibre-wise volume in the free root coordinates]\label{prop:root-fibre-volume}
Fix the Iwasawa diagonal and the middle Levi component.
For \(S,S'\le R_0\) and \(y<Y\le2y\), where \(R_0\ge1\) and \(y>0\),
write \(\mu\) for Lebesgue measure modulo the free roots at these fixed coordinates. The volume estimate \eqref{eq:fibre-volume} and the null-fibre assertion are 
\[
 \begin{aligned}\mu\{S,S'\le R_0,\ y<Y\le2y\}&\le(4r)^{11r^2}R_0^{20r^3}D_e^{12r^2}(\pi e/h)^{rh}(2y)^{2h},\\ \mu\{S,S'\le R_0,\ Y=0\}&=0.\end{aligned}
\]

\end{proposition}

\begin{proof}
Left multiplication by a block-upper unitriangular matrix inside
each block eliminates the within-block entries of the upper
unipotent matrix. The remaining coordinates are the off-block
entries, \(r\) middle columns \(B_b\in\C^{k+1}\), and a
skew-Hermitian \(r\)-by-\(r\) matrix. The within/off-block
factorization has triangular Jacobian one. Relative to a
positive diagonal and a middle Levi component \(m_0\),
the bounds on these coordinates are, respectively,
\[
 R_1=R_0^{2r}D_e,\qquad
 \|m_0^{-1} B_b\|\le M=rR_0^{2r-1}D_e,\qquad
 R_3=rR_0^{4r-2}D_e^2.
\]

\proofstep{Coordinate bounds on the fibre.}

For an off-block row, let \(A_l\) be the leading pairing
matrix at the end \(l\) of its block. The required coordinates
solve \(A_l z=b\) for the pairing vector \(b\) of the remaining
column. If \(A_l[i:b]\) replaces column \(i\) by \(b\), then
the column subtractions in the flag and Hadamard's bound give
\[
 |z_i|=\frac{|\det A_l[i:b]|}{|\det A_l|}
           \le R_0^{2l}D_e\le R_1 .
\]

For a middle column use the \(r\)-by-\(r\) dual isotropic
coordinate matrix \(A\) of the opposite frame.
The expression \(A^{-1}b=(\operatorname{adj}A)b/\det A\)
uses \(r\) cofactors in each coordinate. The orthogonalized
wedge of \(r-1\) columns controls each cofactor. Keeping the
positive diagonal in both this bound and the last column
bound gives \(R_0^{2r-1}D_e\) for each summand and \(M\) after
the sum. With the same inverse, the upper isotropic block
bounds the entries of \(C+H\) by \(R_3\), where \(C\) is
skew-Hermitian and \(H\) is its Hermitian quadratic correction.
The quadratic correction satisfies \(H^*=H\), whereas \(C^*=-C\).
Hence the entry bound passes to
\[
        C=\tfrac12\bigl((C+H)-(C+H)^*\bigr).
\]
\label{cor:absolute-amplitude-2}

In this normal form, the positive projection space is contained
in the middle \(\C^{k+1}\). Its description there is the kernel
of three linear functionals. The negative coordinate vanishes
and the images under \(m_0\) pair to zero against both middle
line vectors. Their \(2\)-by-\(2\) Gram determinant is nonzero
on the cell under consideration. Thus this space has dimension
\(k-2\) or \(k-1\). Let \(\Pi\) be projection onto that middle space.

In the paired
normal form the projected column matrix is
\(\Pi(m_0^{-1}B_b)U\), where \(U\) is the unitriangular
matrix times the opposite isotropic pairing, with the prescribed
adjoints and signs. In particular \(|\det U|=D\). Therefore
\[
  Y=\frac{\|\bigwedge^r(\Pi(m_0^{-1}B_b)U)\|}{D}
    =\|\bigwedge^r(\Pi(m_0^{-1}B_b))\|.
\]
The equality also holds when the projected columns lose rank. The change
\(B_b\mapsto m_0^{-1}B_b\) preserves Lebesgue measure because
\(|\det m_0|=1\).

To bound the volume, project further onto a fixed
\(d=k-2\) dimensional plane in that space. Write \(V\) for
the wedge norm there. Then \(V\le Y\). Integrate
\((2y)^{2h}V^{-2h}\) over the column balls to dominate the
indicator of \(Y\le2y\). Outside the set of linearly dependent
columns, \(V\) is the product of the successive distances
to the previous spans. The exceptional set is null, as is seen
successively by Fubini. Given \(i<r\) previous independent
columns, the integral of the next inverse distance power over
the \(d\)-dimensional complex ball is bounded by
\[
 \frac{\pi^iM^{2i}}{i!}
 \frac{\pi^{d-i} M^{2(d-i-h)}}{(d-i-1)!\,(d-i-h)}
 \le \frac{\pi^d M^{2r}}{h!}.
\]
This follows by integration in polar coordinates perpendicular
to the previous span. Here \(d=r+h\), so \(d-i-h\ge1\).

Integrate the last column first and repeat. The three unused
coordinates per column cost at most \((\pi M^2)^{3r}\).
Using \(h!\ge(h/e)^h\), the middle-column bound is at most
\[
 (\pi M^2)^{3r}M^{2r^2}\pi^{r^2}
                  (\pi e/h)^{rh}(2y)^{2h}.
\]
The other boxes contribute at most
\((\pi R_1^2)^{r(r-1)/2}(\pi R_3^2)^{r(r-1)/2}(2R_3)^r\).
Multiplication, using \(\pi\le4\) and \(r,R_0,D_e\ge1\),
gives \eqref{eq:fibre-volume}. The bound is uniform in the
diagonal, the middle element and the compact right translate.
It also shows that \(Y=0\) is null in the fibres, by letting
the Gram cutoff decrease to zero on each bounded box.

\end{proof}

\subsection{Dyadic summation and the measure factor}

The moment bound controls the decay from one dyadic box to the next. Combining it with the quotient-volume bound on each box brings back the Haar scales in the original orbital integral.

\begin{proposition}[Absolute estimate on a nondegenerate cell]\label{prop:absolute-cell-operator}
For the two frames and the specified quotient measure, assume
\(D=|\Delta_r|>0\), \(r\ge1\), \(k\ge2r+6\),
the moments of Proposition~\ref{prop:paired-frame-moments}
and the Haar scales in \eqref{eq:absolute-root-scales}. With \(h=k-r-2\) there is a constant
\(C\), depending on the fixed family and measure constants, such that
\[
 \|J\|_{\rm op}\le (b^{Cm}D_e)^{Cr^3}
                (C/(kv))^{rh}D^{-h}.
\]
\end{proposition}

\begin{proof}
\proofstep{Shells in the two frame lengths.}

On the \(i'\)-th shell put \(S'\le 2^{i'}H\), with
\(S'>2^{i'-1}H\) if \(i'>0\), and use the analogous shell
\(i\) for \(S\). For \(q=24r^3+1\),
\eqref{eq:side-moment} and the union bound over the \(r+1\)
lengths give relative square mass at most
\((r+1)4^q2^{-2qi'}\) on the primed shell. This covers
\(i'=0\) by the unweighted mass bound. Set
\(R_0=2^{i+i'}H\). Integrating Proposition~\ref{prop:root-fibre-volume} against this
square mass gives, for each unit vector \(w'\),
\begin{equation}\label{eq:side-piece}
 \int_{\substack{i,i'\ {\rm shells}\\y<Y\le2y}}
       |\langle w',\Phi'(x)\rangle|^2\,d\nu_R
 \le K R_0^{20r^3} y^{2h} 2^{-2qi'},
\end{equation}
where
\[
 K=c_N\frac{\kappa_N}{s_R}(4r)^{11r^2}D_e^{12r^2}
       (\pi e/h)^{rh}4^h(r+1)4^q.
\]
Indeed insert a nonnegative quotient cutoff of integral one
over each \(R\)-fibre, and unfold using Weil integration. The
normalized root cutoff for \(d\rho\) then has integral
\(1/s_R\). Integration in \(dN\) gives \(\kappa_N\) times
\eqref{eq:fibre-volume}. The primed shell is invariant under
the primed Bessel subgroup.

On the compact group of the primed side use Haar mass one.
For unit \(w'\), Schur orthogonality gives
\[
 \int_{W'}|\langle w',\pi'(h)\Phi'(x)\rangle|^2\,dh
       = \|\Phi'(x)\|^2/d_{W'} .
\]
The primed shell permits this average inside the quotient integral.
The unweighted square mass is \(c_Nd_{W'}\). The scalar mass
on the primed shell is bounded by
\[
   \frac{1}{d_{W'}}\int_{\{S'\text{ in shell }i'\}}
                   \|\Phi'(x)\|^2d\nu_{S'}(x)
      \le c_N(r+1)4^q2^{-2qi'}.
\]
Here the shell bound uses the moment for the primed frame. The fibre measure contributes its bound with
\(\kappa_N/s_R\), giving \eqref{eq:side-piece}.
The same unfolding transfers the fibrewise null sets to the
quotient. In the reverse orientation send \(x\mapsto\gamma x\)
and \(\rho\mapsto\gamma\rho\gamma^{-1}\).
The transported action is
\[
    \gamma(\rho x)=(\gamma\rho\gamma^{-1})(\gamma x).
\]

It is therefore compatible with quotient integration. To compare
the root measures, consider the reverse normal form.
On a paired block the new upper unitriangular coordinate is
obtained by inverse conjugate transpose and reversal of the
indices, with ratios of pivot coefficients of modulus one.
Inversion on the entry \(u_{ab}\) gives \(-u_{ab}\) plus a
polynomial in entries closer to the diagonal. Its triangular
Jacobian has absolute determinant one. Reindexing, complex
conjugation and multiplication by phases preserve real
Lebesgue measure as well. One obtains the same bound with
\(\Phi(\gamma x),Y',i\) in place of \(\Phi'(x),Y,i'\).

The wedge inequalities bound both positive Gram numbers by
\(T=(rR_0^{2r-1}D_e)^r\), up to the already identified null set.
For integers \(a,b'\ge0\), take boxes
\[
   T2^{-a-1}<Y\le T2^{-a},\qquad T2^{-b'-1}<Y'\le T2^{-b'} .
\]
Denote the two lower endpoints by \(y,y'\). On a nonempty box,
\(yy'\le P=2^r/D\) and \(D\le R_0^{2r}\). Cauchy--Schwarz
on the box and \eqref{eq:side-piece} in both orientations bound
the integral of the product of two absolute scalar coefficients by
\[
 K R_0^{20r^3}(yy')^h2^{-q(i+i')}.
\]
Keep one power of \(yy'\), bounding the others by \(P^{h-1}\).

Since \(1/P\le R_0^{2r}\) and
\(T^2R_0^{2r}=r^{2r}D_e^{2r}R_0^{4r^2}\), the result is at most
\[
 Kr^{2r}D_e^{2r} H^{24r^3}P^h\,
             2^{-i-i'-a-b'} .
\]
We used \(20r^3+4r^2\le24r^3\). All four geometric sums
converge, with product at most 16. A finite basis on each
side gives absolute integrability of the rank-one operator. Taking the supremum over unit test vectors gives the same
bound on the norm of the integral without a dimension loss.
After division by \(c_N\) we obtain
\begin{equation}\label{eq:abs-raw}
 \|J\|_{\rm op}\le \frac{\kappa_N}{s_R}
 (4r)^{44r^3}D_e^{14r^2}H^{24r^3}
        \left(\frac{8\pi e}{k}\right)^{rh}D^{-h}.
\end{equation}
For this simplification use \((k/h)^h\le e^{r+2}\) and
\(4^h(2^r)^h\le8^{rh}\). The remaining factors satisfy
\[
 16(r+1)4^{24r^3+1}r^{2r}e^{r(r+2)}(4r)^{11r^2}
   \le (4r)^{44r^3}
\]
also at \(r=1\).

\proofstep{Root Haar densities.}

In the family normalization put \(N_R=\#\{\text{free complex
roots of }R\}\le r^2\). With local norm parameters
\(\theta,\theta'>0\), the factors in \eqref{eq:abs-raw} are
\begin{equation}\label{eq:absolute-root-scales}
 \kappa_N=(2v)^{-rk}\theta\,2^{r(k+1)+r(r-1)},\qquad
 s_R=\theta'(vg)^{-N_R}.
\end{equation}
For the first identity, the tangent matrix of a unipotent element
contains each complex coordinate and its adjoint partner, and
each imaginary diagonal coordinate once. A paired complex coordinate \(\zeta\) contributes \(2|d\zeta|^2\)
to the tangent metric and gives an area factor of two. The
\(r(k+1)+r(r-1)\) such coordinates give density
\[
        2^{r(k+1)+r(r-1)}(2v)^{-rk}\theta
\]
after the specified Haar scaling.

For the second identity, the root
coordinates of the stabilizer of a normal cell are exactly the
strict upper entries within each block. Reduction by the radical
fixes the first \(r\) flag vectors, so those entries recover the
block coordinates after conjugation. Their permutation preserves
product Lebesgue measure. Each free root has density
\(\kappa_{\rm host}/(\pi g)=1/(vg)\), where
\(\kappa_{\rm host}=\pi/v\). The allowed scalar on the root
measure is \(\theta'\). This also covers a stabilizer with no
free roots.

\proofstep{Insertion of the density ratio.}

For \(b^{-|C_1|m^2}\le\theta,\theta'\le b^{|C_1|m^2}\),
\(1\le v\le b/2\), \(1\le g\le b^{2C_g}\), this gives
\[
 \frac{\kappa_N}{s_R}
 \le b^{2|C_1|m^2+(2+2C_g)r^2} v^{-rh}.
\]
More precisely the formula for the measures gives
\[
  \kappa_N/s_R=2^{r^2}(\theta/\theta')v^{-rk+N_R}g^{N_R}.
\]
Since \(N_R\le r^2\) and \(v\ge1\), the \(v\)-power is at most
\(v^{-rh}\). The remaining factors use \(2\le b\) and the bounds
on \(g,\theta,\theta'\). Inserting this bound and
\eqref{eq:side-moment} into \eqref{eq:abs-raw} yields
\[
 \|J\|_{\rm op}\le (b^{Cm}D_e)^{Cr^3}
                (C/(kv))^{rh}D^{-h}
\]
for \(C\ge8|C_1|+2|C_g|+46+24\max(C_M,0)+8\pi e\).
Here \(4r\le b^m\) and \(m^2\le4mr^3\). This is the
unweighted estimate in Section~\ref{sec:row-sums}. Its
leading-minor iteration there supplies the size weight.

\end{proof}

\subsection{Multiplication by a leading minor}

The frame moment in Proposition~\ref{prop:paired-frame-moments}
estimates an integral on a one-sided quotient. To obtain decay at a
specified pair of frames we first express multiplication by a minor
as an identity of orbital operators \emph{at that pair}. The same
identity will act on both norms used in the row sum.

\begin{auxlemma}[A leading minor acting on an orbital amplitude]
\label{lem:minor-orbital-shift}
Fix the two definite types, characters and all Haar measures of
Proposition~\ref{prop:local-operator-normalization}. Write
\(\mu=\lambda+\nu\), \(p=r+1\), \(n=2r+k+1\), and let \(d_\mu\)
be the formal degree for the fixed ambient Haar measure. Use a nonzero
model on each side and set
\[
 M_\mu=\|J_\mu\|_{\rm cell}^2=C_{\rm cell}/d_\mu,\qquad
 U_\mu=M_\mu^{-1/2}J_\mu .
\]
Take \(1\le l\le p\) and \(\|\nu\|_1\le C_nkm^3\) on the marked
family of Lemma~\ref{lem:marked-paired-moment-range}, sufficiently
late in terms of \(C_n\). At a matched cell in its quotient measure
there are coefficients \(a_I\), independent of its representative
\(\gamma\), such that
\begin{align}
 \Delta_l(\gamma)U_\mu(\gamma)
    &=\sum_{|I|=l}a_I U_{\mu+\mathbf1_I}(\gamma),
       \label{eq:minor-shift-unit}\\
 \Delta_l(\gamma)J_\mu(\gamma)
    &=\sum_{|I|=l}a_I\sqrt{d_{\mu+\mathbf1_I}/d_\mu}\,
                       J_{\mu+\mathbf1_I}(\gamma).
       \label{eq:minor-shift-raw}
\end{align}
A summand with a vanishing period on either side is omitted.
For the choice \(H=b^{m\max(C_M,0)}\) in \eqref{eq:side-moment},
the sum of the absolute coefficients in either line is at most
\((nH)^{2l}\).
\end{auxlemma}

\begin{proof}
Let \(F_l=\bigwedge^l\C^n\), with Hermitian form induced by
\(H_{\rm Witt}\). Let \(\omega,\omega'\) be the ordered wedges of
the first \(l\) columns on the two sides. With the pairing linear
in its second argument take
\[
 \Delta_l(\gamma)=
   [\omega,(\bigwedge\nolimits^l\gamma)\omega']_{F_l}.
\]
This is the determinant of the prefix pairing, including its last
column \(z,z'\) when \(l=p\). Each period subgroup fixes its own
prefix wedge since it adds only earlier columns. Consequently this
finite-dimensional matrix coefficient is constant on the double orbit.

Tensor the ambient representation with \(F_l\). Its weights are the
multiplicity-one vectors \(\mathbf1_I\), \(|I|=l\). The translation
principle in a regular chamber gives, on compact-finite vectors,
\[
 F_l\otimes{\cal H}_\mu
       =\bigoplus_{|I|=l}{\cal H}_{\mu+\mathbf1_I}.
\]
Indeed the coherent family character identity gives one copy of each
translate. They lie in the same chamber by
\eqref{eq:marked-pieri-windows}, and their distinct infinitesimal
characters split the sum
\cite[\S4, (4.1) and the discussion preceding Proposition~4.3]{mpv-translation}.

The invariant tensor form restricts on summand \(I\) to a sign
\(\epsilon_I\) times its Hilbert form after rescaling the inclusion
\(\iota_I\). Write \(P_I\) for its coordinate projection, so
\(P_I\iota_I=1\).

For one of the Schur-normalized periods \(\ell_\mu\), restriction to
the invariant wedge defines a Bessel map on \({\cal H}_\mu\).
Multiplicity one gives the scalar \(\kappa_I\) through
\[
 \ell_{\mu+\mathbf1_I}\big(P_I(\omega\otimes v)\big)
               =\kappa_I\ell_\mu(v).
\]
Use \(\kappa'_I\) for the other side.

Inserting
\(\sum_I\iota_IP_I=1\) into the tensor matrix coefficient gives
a sum of the shifted matrix coefficients with signs \(\epsilon_I\).
Integrating on the original cell is legitimate term by term by the
finite sum and the large-gap majorants in
Proposition~\ref{prop:local-operator-normalization}. Its
Schur formula contributes \(d_I/d_\mu\) to the coefficient of
the shifted standard amplitude, where \(d_I=d_{\mu+\mathbf1_I}\).

If \(J_\mu=\eta_\mu J_{\mu,{\rm std}}\), the full coefficient on
\(J_{\mu+\mathbf1_I}\) and the one on \(U_{\mu+\mathbf1_I}\) are
\[
 b_I=\epsilon_I\kappa'_I\overline{\kappa_I}
       \frac{d_I}{d_\mu}\frac{\eta_\mu}{\eta_{\mu+\mathbf1_I}},
 \qquad
 a_I=b_I\sqrt{d_\mu/d_I}.
\]
Indeed \(|\eta_\mu|=d_\mu\sqrt{d_Wd_{W'}}\) and the compact types
are fixed. Consequently the two coefficient moduli are
\[
  |b_I|=|\kappa_I\kappa'_I|,\qquad
  |a_I|=\sqrt{d_\mu/d_I}\,|\kappa_I\kappa'_I| .
\]
Terms with a zero Bessel map give zero in this computation.

\proofstep{Size of the scalar on one side.}
Normalize the one-sided models as isometries by multiplying
\(\ell_\mu,\ell_{\mu+\mathbf1_I}\) by \(\sqrt{d_\mu},\sqrt{d_I}\).
Write their coefficient functions as \(F_{\mu,v},F_{I,y}\).
Let \(y\) be a unit vector of lowest compact type in the shifted
summand and let \(x_j\) be the components of \(\iota_I y\) in an
orthonormal coordinate basis \(e_j\) of the positive majorant of
\(F_l\).

Since the form on the summand has sign \(\epsilon_I\),
\(\sum_j\|x_j\|^2\ge1\). Its signed adjoint and the quotient
isometry give, with \(h_j(g)\) the \(j\)-th coordinate of
\((\bigwedge^l g^{-1})\omega\) in the dual Hermitian basis,
\[
 |\kappa_I|\sqrt{d_\mu/d_I}\,\|x_j\|^2
  =\left|\int h_j(g)\langle F_{I,y}(g),F_{\mu,x_j}(g)\rangle\,d\nu_S\right|.
\]
To see the projection here, expand \(e_j\otimes x_j\) into its
summands before pairing the coefficient with \(F_{I,y}\). The
other summands are orthogonal in the quotient representation since
they are inequivalent discrete series. The surviving inner product
with \(P_I(e_j\otimes x_j)\), using the Hermitian dual coordinate,
has modulus \(\|x_j\|^2\).

Cauchy--Schwarz and \(\int\|F_{\mu,x_j}\|^2=\|x_j\|^2\) bound
\(|\kappa_I|^2(d_\mu/d_I)\|x_j\|^2\) by
\(\int |h_j|^2\|F_{I,y}\|^2\), including \(x_j=0\).
Sum in \(j\) and evaluate the Euclidean wedge norm by
Gram--Schmidt. H\"older and \eqref{eq:side-moment} give
\[
 |\kappa_I|^2\,d_\mu/d_I
 \le \int\!\Big(\prod_{a=0}^{l-1}l_a(g)^2\Big)\|F_{I,y}(g)\|^2d\nu_S
 \le H^{2l}.
\]
For the last inequality the individual moment of order \(2l\)
follows from the \(2(24r^3+1)\)-moment by H\"older in the
probability measure \(\|F_{I,y}\|^2d\nu_S\). The primed side has
the same bound. Both moments are taken at the shifted parameter.

Finally the formal-degree formula gives \(d_I/d_\mu\le n^l\).
For one raised entry at position \(i\) in decreasing order the
increasing Vandermonde factors cost at most
\(\prod_{j>i}(1+1/(j-i))\le n\); the others decrease. Raise the
entries of \(I\) in increasing order of index, so all intermediate
gaps stay positive. There are at most \(n^l\) possible \(I\).
Together the two one-sided estimates and the count give
\[
 \sum_I|a_I|\le n^{3l/2}H^{2l},\qquad
 \sum_I|b_I|\le n^{2l}H^{2l}.
\]
\end{proof}

\begin{auxlemma}[Iteration at a fixed point of a cell]
\label{lem:minor-orbital-iterate}
Take \(B=b^{C_hm^2}\ge(nH)^{2p}\), and set
\(s=\prod_{l=1}^p\max(1,|\Delta_l(\gamma)|/B)\).
Let \(q\ge0\) be an integer with \(pq\le C_nkm^3\).
For \(V=J\) or \(U\), with either operator or Hilbert--Schmidt norm,
\[
 \|V_\lambda(\gamma)\|
 \le s^{-q/p}\max_{\|\nu\|_1\le pq}\|V_{\lambda+\nu}(\gamma)\|.
\]
Here the maximum ranges over integral shifts with both periods nonzero.
The same inequality holds when the maximum is replaced by any common
unweighted upper bound on those norms.
\end{auxlemma}
\begin{proof}
If every minor has absolute value at most \(B\), the assertion uses
\(\nu=0\). Otherwise choose \(l\) realizing the largest ratio
\(R=|\Delta_l(\gamma)|/B>1\). Iterate \eqref{eq:minor-shift-unit}
or \eqref{eq:minor-shift-raw} with this \emph{same} minor. At depth
\(j\) every path has
\[
 \nu=\mathbf1_{I_1}+\cdots+\mathbf1_{I_j},\qquad
 \|\nu\|_1=jl,\qquad
 \sum_{\rm paths}|\text{product of coefficients}|\le B^j .
\]
The last inequality follows by summing over the children of each
path in Lemma~\ref{lem:minor-orbital-shift}. Each path uses the same
point, quotient measure and compact types. Division by
\(|\Delta_l(\gamma)|^q\) gives \(R^{-q}\), and \(s\le R^p\).
Since \(n\le b^{\max(2,C_u)}\), one may fix
\(C_h\ge\max(C_M,0)+\max(2,C_u)\) before choosing \(C_n\) and the
eventual size threshold.
\end{proof}

The absolute estimate was derived at one real place. Before multiplying it over places one must express it in the same coefficient and Haar normalization at every factor.

\begin{proposition}[Normalized tensor product on an adelic term]\label{prop:tensor-cell-norm}
With matching stabilizer measures in the real factors, the absolute
integral over the stabilizer lattice has the stabilizer covolume
as its measure factor. If \(J_v^{\rm aut}\) denotes the orbital operator at \(v\), the squared Hilbert--Schmidt norm including its finite factor is the expression of \eqref{eq:adelic-tensor-scale}: 
\[
 (s_{u,u'}\operatorname{vol}K_2)^2\operatorname{covol}(R_{\rm arith})^2\prod_v\|J_v^{\rm aut}\|_{\rm HS}^2.
\]
 Under the coefficient isometries
the local squared norms gain the squares of the two vector
scalars divided by the square of the Bessel normalization.
\end{proposition}

\begin{proof}
\proofstep{Agreement of stabilizer measures.}

For the application to an adelic term choose the stabilizer
Haar measure by its specified pushforward in the rational
root coordinates. At the real places rescale each local Haar
measure so its free-coordinate pushforward is Lebesgue with
density \((vg)^{-N_R}\). Their product and the pinned global
archimedean Haar measure are proportional, by uniqueness of
Haar measure on the product stabilizer. Their coordinate
pushforwards agree. Testing on a bounded coordinate box of
positive measure fixes the proportionality at one. Left and
right changes of double-coset representative conjugate the
stabilizer on the corresponding side. Normal-form coordinates
and the quotient integration formula transport together under
these maps. Relevance gives equality of the two characters
on the stabilizer, so their phases cancel in the rank-one
integrand. Compact changes of representatives act by unitaries
on the two coefficient spaces.

\proofstep{Tensor factorization and coefficient normalization.}

Set \(c_R=\operatorname{covol}(R_{\rm arith})\). The quotient
integration formula and Fubini give
\[
   \int_{R_{\rm arith}\backslash G_\infty}(\text{paired integrand})
        =c_R\bigotimes_v J_v^{\rm aut},\qquad
   \left\|c_R\bigotimes_v J_v^{\rm aut}\right\|_{\rm HS}^2
        =c_R^2\prod_v\|J_v^{\rm aut}\|_{\rm HS}^2 .
\]
The root measures on each \(J_v^{\rm aut}\) are the ones just
compared. Write \(s_{u,u'}\) for the
product normalization of the two finite-level column systems
and \(R_{\rm arith}\) for their stabilizer lattice. Before
transferring from the automorphic to the Bessel coefficient
normalization, the squared norm of the normalized lift is
\begin{equation}\label{eq:adelic-tensor-scale}
 (s_{u,u'}\operatorname{vol}K_2)^2
      \operatorname{covol}(R_{\rm arith})^2
      \prod_v\|J_v^{\rm aut}\|_{\rm HS}^2 .
\end{equation}
The finite level character has modulus
one.

The local comparison of the two realizations multiplies
\(\|J_v^{\rm aut}\|_{\rm HS}^2\) by the squared modulus of
the product of the two vector-normalization scalars divided by
the squared Bessel normalization. This follows by substituting
the two transformed coefficient functions in the rank-one
integral. It applies at a Witt lift of the term because the
transformed quotient measure is the one for its transformed
root measure. This explains which covolume and which local
norms can occur together in the row product.
\label{cor:absolute-amplitude-3}
\end{proof}
